%=====================================================================
\chapter{Datasets and Further Resources}\label{app:resources}
%=====================================================================
\normalfigures

\section{Datasets for Practice}

Grouped by the four running domains, so that you can practise a chapter's technique
in the lens that interests you.

\rowcolors{2}{white}{lightbg}
\begin{longtable}{@{}L{4.2cm}L{3.0cm}L{6.8cm}@{}}
\toprule
\rowcolor{primary!15}
\textbf{Dataset} & \textbf{Domain} & \textbf{Good for practising} \\
\midrule
\endhead
Open University Learning Analytics (OULAD) & Learning analytics & The whole of this course: clickstream, demographics, outcomes, and genuine class imbalance (Ch 7, 15, 19) \\
Student Performance (UCI) & Learning analytics & Small and clean; ideal for Chapters 5--12 \\
xAPI Educational Data & Learning analytics & Engagement features and fairness slices (Ch 28) \\
\addlinespace
GitHub / JIRA public issue exports & Project management & Small-$n$ modelling, effort estimation, the honest limits of Ch 30 \\
PROMISE software engineering repository & Project management & Defect prediction with realistic sample sizes \\
\addlinespace
NASA C-MAPSS turbofan degradation & IoT & Remaining useful life, sensor windows, the whole of Ch 25 \\
UCI Air Quality / SECOM & IoT & Missing sensor data, drift, high-dimensional telemetry \\
Intel Berkeley Lab sensor data & IoT & Raw telemetry with genuine dropouts and clock issues \\
\addlinespace
CICIDS2017 / CSE-CIC-IDS2018 & Cybersecurity & Intrusion detection with realistic class imbalance (Ch 14, 26) \\
UNSW-NB15 & Cybersecurity & Modern attack categories; good for PR-curve work \\
Enron email corpus & Cybersecurity / NLP & Phishing-style text classification (Ch 20) \\
EMBER (malware features) & Cybersecurity & Malware classification without handling binaries \\
\addlinespace
\texttt{mtcars}, \texttt{palmerpenguins}, \texttt{titanic} & Teaching & Small, clean, well documented; the department's existing R labs use \texttt{mtcars} \\
\bottomrule
\end{longtable}

\begin{alertbox}[Before You Download Anything]
Check the licence and the permitted use, and check whether the data contains
personal information (\chref{ethics}). Several widely-circulated ``teaching''
datasets contain real personal data released without adequate consent, and using
them uncritically is not a neutral act. If a dataset has no documented provenance
(\chref{data}), treat any result from it as provisional.
\end{alertbox}

\section{Where to Find Data}

\begin{tight}
  \item \textbf{UCI Machine Learning Repository} --- the classic teaching source;
        small, clean, well documented.
  \item \textbf{Kaggle Datasets} --- large and varied; quality varies just as much,
        so read the provenance notes.
  \item \textbf{Government open data portals} --- real, messy, and usually well
        licensed. Excellent for \chref{wrangling}.
  \item \textbf{Your own institution} --- with permission, the most valuable source
        for a capstone, because the results can actually be used.
  \item \textbf{Papers With Code} --- datasets attached to published benchmarks.
\end{tight}

\section{Tools}

\rowcolors{2}{white}{defbg}
\begin{longtable}{@{}L{3.4cm}L{4.2cm}L{6.4cm}@{}}
\toprule
\rowcolor{primary!15}
\textbf{Purpose} & \textbf{Tool} & \textbf{Note} \\
\midrule
\endhead
Analysis (Python) & pandas, NumPy, scikit-learn & The core stack of this handout \\
Analysis (R) & tidyverse, tidymodels & Better for exploratory statistics and regression diagnostics \\
Notebooks & Jupyter, Quarto & Quarto renders both R and Python and produces the report \\
Visualisation & matplotlib, seaborn, ggplot2 & \chref{eda} \\
Gradient boosting & XGBoost, LightGBM, CatBoost & \chref{features} \\
Deep learning & PyTorch, TensorFlow/Keras & \chref{neuralnets} \\
Time series & statsmodels, Prophet, sktime & \chref{timeseries} \\
Text & scikit-learn, spaCy, sentence-transformers & \chref{nlp} \\
Explainability & SHAP, LIME, \texttt{sklearn.inspection} & \chref{ethics} \\
Fairness & Fairlearn, AIF360 & Subgroup metrics; still requires your judgement \\
Pipelines & Airflow, Dagster, dbt & \chref{dataeng} \\
Experiment tracking & MLflow, Weights \& Biases & \chref{mlops} \\
Containers & Docker & \chref{cloudedge} \\
Version control & git, DVC (for data) & \chref{research} \\
\bottomrule
\end{longtable}

\section{Reading Beyond This Handout}

Grouped by what you would go to them \emph{for}, rather than alphabetically.

\begin{tight}
  \item \textbf{To go deeper on the statistics of Part II} --- \emph{Statistical
        Rethinking} (McElreath) for a modern Bayesian treatment;
        \emph{Introduction to Statistical Learning} (James et al.) for the
        statistics of Part III, free and with both R and Python editions.
  \item \textbf{To go deeper on Part III} --- \emph{Hands-On Machine Learning}
        (Géron) is the standard practical companion; \emph{The Elements of
        Statistical Learning} (Hastie et al.) is the theoretical reference.
  \item \textbf{To go deeper on Part IV} --- \emph{Deep Learning} (Goodfellow et
        al.) for theory; the fast.ai course for practice.
  \item \textbf{On causality (\chref{causality})} --- \emph{The Book of Why}
        (Pearl) for intuition; \emph{Causal Inference: The Mixtape} (Cunningham)
        for the quasi-experimental designs.
  \item \textbf{On engineering (Part V)} --- \emph{Designing Data-Intensive
        Applications} (Kleppmann); \emph{Designing Machine Learning Systems}
        (Huyen) covers \chref{mlops} particularly well.
  \item \textbf{On ethics (\chref{ethics})} --- \emph{Weapons of Math Destruction}
        (O'Neil) for the argument; \emph{Fairness and Machine Learning}
        (Barocas, Hardt, Narayanan) for the technical treatment, free online.
  \item \textbf{On communication (\chref{communication})} --- \emph{Storytelling
        with Data} (Knaflic); \emph{The Visual Display of Quantitative
        Information} (Tufte).
  \item \textbf{On research method (\chref{research})} --- your institution's own
        project handbook takes precedence over any general text, because it is what
        you are marked against.
\end{tight}

\section{Where the Spring 2026 Material Went}

This volume consolidates the thirteen separate handouts used in the previous
offering. If you are looking for one of them, this is where its content now lives.

\rowcolors{2}{white}{lightbg}
\begin{longtable}{@{}L{6.4cm}L{7.6cm}@{}}
\toprule
\rowcolor{primary!15}
\textbf{Spring 2026 handout} & \textbf{Now covered by} \\
\midrule
\endhead
Foundations of Data Science, AI and Emerging Paradigms & Chapters 1--2, with the AI material distributed to 17--22 \\
Mathematical Foundations of Data Science & Chapter 4 \\
Statistical Foundations of Data Science & Chapters 5, 6, 9 \\
Machine Learning & Chapters 11--16 \\
Deep Learning & Chapters 17--18 \\
Advanced Data Science & Chapters 23, 27, 28 \\
Generative AI & Chapter 21 \\
Agentic AI / AI Agents & Chapter 22 \\
Cloud Computing & Chapter 24 \\
IoT for Data Science & Chapter 25 \\
Cybersecurity & Chapter 26 \\
IT Project Management with AI & Chapter 30 \\
Research Methods in IT & Chapter 29 \\
\bottomrule
\end{longtable}

\begin{conceptbox}[What Changed Besides the Order]
Three things. First, material that appeared in several handouts --- the ML/DL
overview in particular --- now appears exactly once. Second, every chapter carries a
Four Lenses panel, so the application domains are present throughout rather than
confined to their own handouts. Third, IoT and cybersecurity are now positioned
after the methods they depend on, so they can be taught properly rather than
described.
\end{conceptbox}
